\section{Related Work}
\label{sec:related}

\subsection{Reinforcement Learning from Execution Feedback for Code LLMs}

The dominant paradigm for RLEF post-training of code LLMs uses binary pass/fail reward: the model receives reward 1 if and only if all test cases pass, and 0 otherwise. CodeRL \citep{le2022coderl} pioneered this approach with PPO on the APPS dataset \citep{hendrycks2021apps}, reporting approximately 5 percentage-point improvements over supervised fine-tuning alone on HumanEval. PPOCoder \citep{shojaee2023ppocoder} extended this framework with additional reward components including syntax correctness signals, but retained binary execution reward as the primary training signal. RLEF \citep{gehring2024rlef} applied binary RLEF to repository-level code repair on SWE-bench, demonstrating that execution-grounded reward can produce gains over SFT even in complex agentic settings.

All three works treat reward formulation as a fixed design choice. None analyze whether binary reward's group-level gradient properties are optimal for GRPO, nor do they compare binary reward against ratio or continuous alternatives in controlled ablation. Our work asks a question these papers did not ask: under what conditions does binary reward produce zero gradient contribution in GRPO training, and does ratio reward address this?

\subsection{Group Relative Policy Optimization}

GRPO \citep{shao2024deepseekmath} computes advantages by normalizing rewards within a group of completions, eliminating the need for a separate critic network. This simplification reduces training complexity and memory but changes the reward sparsity problem qualitatively: where PPO's value function provides a baseline even when execution rewards are sparse, GRPO's group normalization can only redistribute signal within a group --- it cannot create signal where none exists. If all completions in a group receive the same reward (including all-zero under binary reward), GRPO produces zero gradient contribution from that group.

DAPO \citep{yu2025dapo} acknowledges the sparsity concern and introduces several mitigations (dynamic sampling temperature, clip-higher variant). However, DAPO retains binary reward and does not quantify the fraction of training groups that fall into the all-zero regime. Our analysis shows that under realistic early-training conditions ($p_{\text{pass}} \approx 0.1$ per test case), 98.7\% of GRPO groups receive zero binary reward for all completions --- a dead zone that DAPO's mitigations do not directly address for hard coding tasks.

\subsection{Reward Signal Density and Sparsity}

The problem of sparse rewards in reinforcement learning has a long history \citep{sutton2018rl}. Reward shaping \citep{ng1999policy} addresses sparsity at the policy level by adding potential-based bonuses that preserve the optimal policy. The ratio reward we analyze ($k/n$ test cases passing) is an instance of partial-credit reward --- a natural generalization of binary reward. Our contribution is not the reward function itself, but the group-level analysis showing when it provides a mathematical guarantee of gradient signal that binary reward cannot provide, and the precise characterization of when this guarantee is active in practice.

\subsection{Positioning}

Our work differs from all prior RLEF literature by analyzing gradient signal quality at the GRPO group level --- the atomic unit of computation --- rather than at the aggregate training curve level. This framing reveals the binary reward dead zone that is invisible in aggregate metrics, and provides a precise, falsifiable condition for when ratio reward outperforms binary reward in terms of gradient signal density.
